%======================================================================
\section{Background and the setup}
\label{sec:setup}
%======================================================================

\subsection{Light-cone coordinates and conventions}
\label{sec:conv}

We use the light-cone conventions of Paper~I \cite{PaperI,Harindranath:1996hq,Burkardt:1995ct}:
$x^{\pm}=(x^{0}\pm x^{3})/\sqrt2$, quantization on surfaces of fixed $x^{+}$, longitudinal momentum
$p^{+}>0$ for physical particles, $d^{3}p\equiv dp^{+}d^{2}\bm p$ and
$\delta^{(3)}(k-p)\equiv\delta(k^{+}-p^{+})\delta^{(2)}(\bm k-\bm p)$. For an on-shell massless
particle the light-cone energy is $p^{-}=\bm p^{2}/2p^{+}$, which we abbreviate as
\begin{equation}
\eps_{p}\equiv\frac{\bm p^{2}}{2p^{+}} .
\label{eq:eps}
\end{equation}
We also use the shorthand $\int_{p}\equiv\int d^{3}p/(2\pi)^{3}$. We work in light-cone gauge
$A^{+}=0$, in which only the transverse gluon fields $A^{a}_{i}$, $i=1,2$, are dynamical; their
creation and annihilation operators satisfy
\begin{equation}
\big[a^{a}_{i}(k),a^{\dg b}_{j}(p)\big]=(2\pi)^{3}\delta^{ab}\delta_{ij}\,\dthree(k-p).
\label{eq:aa}
\end{equation}
The quark field is decomposed as $\psi=\psi_{+}+\psi_{-}$ with $\psi_{\pm}=\Lambda_{\pm}\psi$,
$\Lambda_{\pm}=\tfrac12\gamma^{\mp}\gamma^{\pm}$. Only the ``good'' component $\psi_{+}$ is
dynamical; $\psi_{-}$ is constrained and is eliminated through its equation of motion
(appendix~\ref{app:hamiltonian}). In the two-component representation
\cite{Zhang:1993dd,Brodsky:1997de} $\psi_{+}$ is described by a two-component field $\xi$, which we
expand as
\begin{equation}
\xi^{\beta}(x)=2^{-1/4}\!\int_{0}^{\infty}\!\frac{dk^{+}}{2\pi}\!\int\!\frac{d^{2}k}{(2\pi)^{2}}
\sum_{\lambda=\pm}\Big[b^{\beta}_{\lambda}(k)\,\chi_{\lambda}\,e^{-ik\cdot x}
+d^{\dg\beta}_{\lambda}(k)\,\chi_{\lambda}\,e^{ik\cdot x}\Big],
\label{eq:xi}
\end{equation}
where $\beta$ is a colour index in the fundamental representation, $\chi_{\lambda}$ are the
eigenvectors of $\sigma^{3}$ with $\sum_{\lambda}\chi_{\lambda}\chi^{\dg}_{\lambda}=\mathbf 1$, and a
flavour index is suppressed.\footnote{With this labelling the physical helicity of an antiquark
created by $d^{\dg}_{\lambda}$ is $-\lambda$. It has the advantage that every quark--gluon vertex
reduces to a single spinor bilinear of the form \eqref{eq:Phi}.} The quark and antiquark
operators obey
\begin{equation}
\big\{b^{\beta}_{\lambda}(l),b^{\dg\beta'}_{\lambda'}(l')\big\}
=\big\{d^{\beta}_{\lambda}(l),d^{\dg\beta'}_{\lambda'}(l')\big\}
=(2\pi)^{3}\delta^{\beta\beta'}\delta_{\lambda\lambda'}\,\dthree(l-l'),
\label{eq:bb}
\end{equation}
all other anticommutators vanishing; fermion operators commute with gluon operators.

All the quark--gluon vertices of light-cone QCD follow from one structure. For a gluon of
momentum $k$ and transverse polarization $i$ attached to a fermion line, let $P_{L}$ and $P_{R}$
be the momenta of the two fermion legs, $P_{L}$ being the one at the head of the fermion number
arrow and $P_{R}$ the one at its tail (antiquark momenta are physical, i.e.\ positive, momenta).
We define
\begin{equation}
\Gam^{i}(k;P_{L},P_{R})\equiv\frac{2k^{i}}{k^{+}}
-\frac{(\bm\sigma\!\cdot\!\bm P_{L})\,\sigma^{i}}{P_{L}^{+}}
-\frac{\sigma^{i}\,(\bm\sigma\!\cdot\!\bm P_{R})}{P_{R}^{+}},
\qquad
\Ph^{i}_{\lambda\lambda'}(k;P_{L},P_{R})\equiv\chi^{\dg}_{\lambda}\,\Gam^{i}(k;P_{L},P_{R})\,\chi_{\lambda'} .
\label{eq:Phi}
\end{equation}
The three elementary processes then carry the assignments
\begin{center}
\begin{tabular}{lcc}
\toprule
process & $P_{L}$ & $P_{R}$\\\midrule
$g(k)\to q(l)\,\bar q(m)$ & $l$ (quark) & $m$ (antiquark)\\
$q(l')\to q(l)\,g(k)$ & $l$ (outgoing quark) & $l'$ (incoming quark)\\
$\bar q(m')\to\bar q(m)\,g(k)$ & $m'$ (incoming antiquark) & $m$ (outgoing antiquark)\\
\bottomrule
\end{tabular}
\end{center}
and the conjugate (absorption) processes follow from $\big[\Gam^{i}(k;P_{L},P_{R})\big]^{\dg}
=\Gam^{i}(k;P_{R},P_{L})$. Since
$2k^{i}-(\bm\sigma\!\cdot\!\bm k)\sigma^{i}-\sigma^{i}(\bm\sigma\!\cdot\!\bm k)=0$, the vertex
depends only on the relative transverse momentum of the splitting. For $g(k)\to q(l)\bar q(m)$,
with $z=l^{+}/k^{+}$ and $\bm\kappa=\bm l-z\bm k$,
\begin{equation}
\Gam^{i}(k;l,m)=\frac{(2z-1)\kappa^{i}+i\epsilon^{ij}\kappa^{j}\sigma^{3}}{k^{+}z(1-z)},
\qquad
\eps_{k}-\eps_{l}-\eps_{m}=-\frac{\bm\kappa^{2}}{2z(1-z)k^{+}} ,
\label{eq:kappa}
\end{equation}
and for $q(l')\to q(l)g(p)$, with $z=p^{+}/l'^{+}$ and $\bm\kappa=\bm p-z\bm l'$,
\begin{equation}
\Gam^{j}(p;l,l')=\frac{(2-z)\kappa^{j}-iz\,\epsilon^{jk}\kappa^{k}\sigma^{3}}{l'^{+}z(1-z)},
\qquad
\eps_{l'}-\eps_{l}-\eps_{p}=-\frac{\bm\kappa^{2}}{2z(1-z)l'^{+}} .
\label{eq:kappaq}
\end{equation}
Finally, the three-gluon vertex of Paper~I is written as
\begin{equation}
\Vg^{ijl}(k;p,q)=\Big(p^{i}-q^{i}+\tfrac{q^{+}-p^{+}}{k^{+}}k^{i}\Big)\delta^{jl}
+\Big(k^{j}+q^{j}-\tfrac{k^{+}+q^{+}}{p^{+}}p^{j}\Big)\delta^{il}
+\Big(\tfrac{k^{+}+p^{+}}{q^{+}}q^{l}-p^{l}-k^{l}\Big)\delta^{ij}.
\label{eq:Vggg}
\end{equation}

\subsection{Light-cone QCD Hamiltonian with quarks}
\label{sec:ham}

The derivation of the QCD Hamiltonian in light-cone gauge is standard
\cite{Brodsky:1997de,Zhang:1993dd,Kogut:1969xa,Lublinsky:2016meo}; it is summarized in
appendix~\ref{app:hamiltonian}. After the constrained fields $A^{-}$ and $\psi_{-}$ are eliminated,
the Hamiltonian reads
\begin{equation}
H_{\rm LC\,QCD}=H_{\rm LC\,YM}+H_{q},
\label{eq:HQCD}
\end{equation}
where $H_{\rm LC\,YM}$ is the light-cone Yang--Mills Hamiltonian of Paper~I and, for massless
quarks,
\begin{align}
H_{q}=\int\!dx^{-}d^{2}x\,\Big[\;&\frac{1}{2}\bar\psi\gamma^{+}\frac{-\bm\partial^{2}}{i\partial^{+}}\psi
+g\,\bar\psi\gamma^{i}t^{a}\psi\,A^{a}_{i}
+\frac{g^{2}}{2}\,\bar\psi\gamma^{i}t^{a}A^{a}_{i}\,\frac{\gamma^{+}}{2i\partial^{+}}\,
\gamma^{j}t^{b}A^{b}_{j}\,\psi
\nonumber\\
&+g^{2}\,J^{+a}_{q}\frac{1}{(i\partial^{+})^{2}}J^{+a}_{g}
+\frac{g^{2}}{2}\,J^{+a}_{q}\frac{1}{(i\partial^{+})^{2}}J^{+a}_{q}\;\Big],
\label{eq:Hq}
\end{align}
with the quark and gluon light-cone colour currents
\begin{equation}
J^{+a}_{q}=\bar\psi\gamma^{+}t^{a}\psi,\qquad J^{+a}_{g}=f^{abc}A^{b}_{i}\partial^{+}A^{c}_{i}.
\label{eq:currents}
\end{equation}
The first term of \eqref{eq:Hq} is the free quark Hamiltonian; the second is the quark--gluon
vertex; the third is the instantaneous quark exchange (it arises from eliminating $\psi_{-}$);
the last two are the instantaneous gluon exchange terms involving the quark current (they
arise, together with the $J_{g}J_{g}$ term contained in $H_{\rm LC\,YM}$, from eliminating
$A^{-}$). The instantaneous gluon exchange is most usefully regarded as a single current--current
interaction,
\begin{equation}
H^{(A)}_{\rm inst}=\frac{g^{2}}{2}\int\!dx^{-}d^{2}x\;\Jc^{a}\frac{1}{(i\partial^{+})^{2}}\Jc^{a},
\qquad \Jc^{a}=J^{+a}_{\rm val}+J^{+a}_{g}+J^{+a}_{q},
\label{eq:HinstA}
\end{equation}
because after the Born--Oppenheimer separation discussed below the classification of its
contributions becomes automatic: one expands $\Jc$ and reads off the cross terms. The free
Hamiltonian of the soft sector is
\begin{equation}
H_{0}=\int_{k}\eps_{k}\,a^{\dg a}_{i}(k)a^{a}_{i}(k)
+\int_{l}\eps_{l}\sum_{\lambda,\beta}\Big[b^{\dg\beta}_{\lambda}(l)b^{\beta}_{\lambda}(l)
+d^{\dg\beta}_{\lambda}(l)d^{\beta}_{\lambda}(l)\Big],
\label{eq:H0}
\end{equation}
whose eigenstates are the free Fock states of soft gluons, quarks and antiquarks.

\subsection{Born--Oppenheimer formalism and the eikonal approximation}
\label{sec:BO}

As in Paper~I, the modes with longitudinal momentum $k^{+}\ge\vee$ are the valence degrees of
freedom and those with $\Lambda<k^{+}<\vee$ the soft ones \cite{McLerran:1993ni,McLerran:1993ka}.
With quarks, both the gluon and the quark fields are decomposed into valence and soft
components,
\begin{equation}
A^{a}_{i}=\mathcal A^{a}_{i}+A^{a}_{i},\qquad
\xi=\Xi+\xi ,
\label{eq:split}
\end{equation}
where calligraphic (capital) symbols denote the valence modes. Substituting into the Hamiltonian
generates soft--soft, soft--valence and valence--valence interactions. Since the valence modes
carry parametrically larger longitudinal momenta, contributions suppressed by powers of
$k^{+}_{\rm soft}/k^{+}_{\rm valence}$ are neglected (eikonal approximation), as are terms in which
inverse longitudinal derivatives act on valence fields \cite{Lublinsky:2016meo,PaperI}. In
particular, a soft quark cannot absorb a valence gluon (nor a soft gluon a valence quark) without
leaving the soft sector, and the transverse couplings $\bar\psi\gamma^{i}\psi A_{i}$ of the
valence quarks are power suppressed. The only coupling of the valence sector to the soft modes
that survives is through the plus component of the valence colour current, i.e.\ through the
total valence colour charge density
\begin{equation}
\rho^{a}(-p)=\rho^{a}_{g}(-p)+\rho^{a}_{q}(-p),
\label{eq:rho}
\end{equation}
with the gluonic part of Paper~I and a fundamental representation part carried by the valence
quarks and antiquarks,
\begin{align}
\rho^{a}_{g}(-p)&=-if^{abc}\!\int_{\vee}^{\infty}\!\frac{dk^{+}}{2\pi}\!\int\!\frac{d^{2}k}{(2\pi)^{2}}\;
\bar a^{\dg b}_{j}(k)\,\bar a^{c}_{j}(k+p),
\label{eq:rhog}\\
\rho^{a}_{q}(-p)&=\int_{\vee}^{\infty}\!\frac{dk^{+}}{2\pi}\!\int\!\frac{d^{2}k}{(2\pi)^{2}}
\sum_{\lambda}\Big[\bar b^{\dg\alpha}_{\lambda}(k)\,t^{a}_{\alpha\beta}\,\bar b^{\beta}_{\lambda}(k+p)
-\bar d^{\dg\alpha}_{\lambda}(k+p)\,t^{a}_{\beta\alpha}\,\bar d^{\beta}_{\lambda}(k)\Big],
\label{eq:rhoq}
\end{align}
barred operators denoting valence modes and a sum over flavours being implied. Both pieces
satisfy the same algebra, so that
\begin{equation}
\big[\rho^{a}(k),\rho^{b}(p)\big]=if^{abc}\rho^{c}(k+p)
\label{eq:rhoalg}
\end{equation}
holds for the total charge density. Every result of Paper~I therefore carries over verbatim
with the enlarged $\rho$; what changes is the operator content of $\rho$ and, more importantly,
the appearance of soft quark degrees of freedom.

After performing the mode decomposition and applying the eikonal approximation, the soft
interaction Hamiltonian of QCD can be organized as
\begin{equation}
H_{\rm int}=\underbrace{H_{g}+H_{ggg}+H_{gggg}+H_{gg-{\rm inst}}+H_{gggg-{\rm inst}}}_{\text{Paper~I}}
+H_{gq\bar q}+H_{ggq\bar q}+H_{q-{\rm inst}}+H_{V},
\label{eq:Hint}
\end{equation}
where the first five terms are those of Paper~I with $\rho$ given by \eqref{eq:rho}, and $H_{V}$
contains valence--valence interactions only. In particular
\begin{equation}
H_{g}=\int_{k}\frac{g\,k^{i}}{\sqrt2\,|k^{+}|^{3/2}}
\Big[a^{\dg a}_{i}(k)\rho^{a}(-k)+a^{a}_{i}(k)\rho^{a}(k)\Big].
\label{eq:Hg}
\end{equation}
The quark--gluon vertex of the soft sector, written in terms of the vertex function
\eqref{eq:Phi}, contains three structures and their Hermitian conjugates,
\begin{align}
H_{gq\bar q}=\int_{k,l,m}&\Big[V^{a\alpha\beta}_{3\,i\lambda_{2}\lambda_{1}}(k,m,l)\,
b^{\dg\beta}_{\lambda_{1}}(l)d^{\dg\alpha}_{\lambda_{2}}(m)a^{a}_{i}(k)+{\rm h.c.}\Big]
\nonumber\\
+\int_{p,l,l'}&\Big[V^{\beta'\beta b}_{F\,\lambda'\lambda j}(l',l,p)\,
a^{\dg b}_{j}(p)b^{\dg\beta}_{\lambda}(l)b^{\beta'}_{\lambda'}(l')+{\rm h.c.}\Big]
\nonumber\\
+\int_{p,m,m'}&\Big[V^{\alpha'\alpha b}_{G\,\lambda'\lambda j}(m',m,p)\,
a^{\dg b}_{j}(p)d^{\dg\alpha}_{\lambda}(m)d^{\alpha'}_{\lambda'}(m')+{\rm h.c.}\Big],
\label{eq:Hqqg}
\end{align}
with
\begin{align}
V^{a\alpha\beta}_{3\,i\lambda_{2}\lambda_{1}}(k,m,l)&=\frac{(2\pi)^{3}g\,t^{a}_{\beta\alpha}}{2\sqrt2\sqrt{k^{+}}}\,
\dthree(k-l-m)\,\Ph^{i}_{\lambda_{1}\lambda_{2}}(k;l,m),
\label{eq:V3}\\
V^{\beta'\beta b}_{F\,\lambda'\lambda j}(l',l,p)&=\frac{(2\pi)^{3}g\,t^{b}_{\beta\beta'}}{2\sqrt2\sqrt{p^{+}}}\,
\dthree(l'-l-p)\,\Ph^{j}_{\lambda\lambda'}(p;l,l'),
\label{eq:VF}\\
V^{\alpha'\alpha b}_{G\,\lambda'\lambda j}(m',m,p)&=-\frac{(2\pi)^{3}g\,t^{b}_{\alpha'\alpha}}{2\sqrt2\sqrt{p^{+}}}\,
\dthree(m'-m-p)\,\Ph^{j}_{\lambda'\lambda}(p;m',m).
\label{eq:VG}
\end{align}
The first line of \eqref{eq:Hqqg} creates (annihilates) a soft quark--antiquark pair from (into) a
soft gluon; the second and third describe the emission and absorption of a soft gluon by a soft
quark and by a soft antiquark. The minus sign and the transposed colour matrix in \eqref{eq:VG}
reflect the antiquark charge. The instantaneous quark exchange of \eqref{eq:Hq} restricted to soft
fields is denoted by $H_{ggq\bar q}$. Expanding the current $\Jc$ of \eqref{eq:HinstA}, the
cross terms $J_{\rm val}\cdot J_{g}$ and $J_{g}\cdot J_{g}$ are the $H_{gg-{\rm inst}}$ and
$H_{gggg-{\rm inst}}$ of Paper~I, while the quark current gives
\begin{equation}
H_{q-{\rm inst}}=g^{2}\!\int\!dx^{-}d^{2}x\;\Big[\rho^{a}\,\frac{1}{(i\partial^{+})^{2}}\,J^{+a}_{q}
+J^{+a}_{g}\frac{1}{(i\partial^{+})^{2}}J^{+a}_{q}
+\frac12\,J^{+a}_{q}\frac{1}{(i\partial^{+})^{2}}J^{+a}_{q}\Big],
\label{eq:Hqinst}
\end{equation}
with all currents built from soft fields except $\rho$. The first term describes the
instantaneous interaction of soft quarks with the valence charge (eikonal colour rotation and
pair creation off $\rho$), the second the instantaneous interaction between the soft gluon and
quark currents, and the third the instantaneous interaction among soft quarks. The matrix
elements of all these operators between soft Fock states are collected in
appendix~\ref{app:me}. We count $\rho=\order{1}$ in the dilute regime, as in Paper~I, so that
$H_{g}$, $H_{ggg}$ and $H_{gq\bar q}$ are of order $g$ and $H_{gggg}$, the instantaneous terms and
$H_{ggq\bar q}$ of order $g^{2}$.
